\section{Background}
This section provides the theoretical foundation for this work while excluding Graph-based representations of constraint problems as we will discuss them in Section \ref{sec:conversion}.

\subsection{Combinatorial Optimisation}
Constraint Programming (CP) \cite{rossi2006handbook} is a paradigm for solving combinatorial optimization problems (COPs) by defining them via variables, domains, and constraints. Unlike other combinatorial solvers, CP utilizes global constraints (such as \textit{all-different}, \textit{cumulative}, or \textit{circuit}) which capture complex substructures. These constraints use specialized filtering algorithms to prune the search space through propagation, enabling more efficient modeling and solving.

Constraint solvers navigate the solution space using techniques like constraint propagation \cite{bessiere2006constraint}, branch and bound \cite{boyd2007branch}, and conflict-driven clause learning (CDCL) \cite{marques2009conflict}. Given that COPs are often NP-hard \cite{Papadimitriou_book}, parallelization is frequently employed to speed up the computation. This typically follows two patterns: (i) the job-splitting approach, which parallelizes a single solver's strategy, and (ii) the portfolio approach, which runs multiple solvers with different strategies simultaneously \cite{paper_amai}. Comprehensive reviews for parallelizing these specific techniques are available in \cite{gent2018review, crainic2006parallel, martins2012overview}.

\subsection{Algorithm Selection}
Algorithm Selection (AS) involves identifying the most suitable algorithm for a specific problem instance to optimize metrics like computational efficiency or accuracy. This process is necessitated by the No Free Lunch theorem \cite{wolpert1997no}, which states that no single algorithm is universally superior. In the SAT domain, SATZilla \cite{xu2012satzilla2012} uses AS to manage solver portfolios. In CP, approaches like CPHydra \cite{bridge2012case} and SUNNY \cite{amadini2014sunny} utilize $k$-nearest neighbors (k-NN) to generate instance-specific solver schedules.

\subsection{Feature Extraction on Graphs}
Feature extraction transforms graph-structured data into fixed-dimensional vector spaces \cite{graphkernels}. Classic methods rely on graph-theoretical metrics to describe structural properties.

Node-level features describe a node's topological importance. Common metrics include Degree Centrality, Clustering Coefficient \cite{watts1998collective}, and Betweenness or Closeness Centrality \cite{newman2010network}. Edge-level features are primarily used for link prediction and node proximity. These include Common Neighbors, the Jaccard Coefficient, the Adamic-Adar Index (which weights neighbors by their rarity) \cite{adamic2003friends}, and the Katz Index \cite{Luce_Perry_1949}.

Graph-level features are often extracted via kernels that compute similarity as a scalar product between two graphs $G_1$ and $G_2$. 
\begin{itemize}
    \item \textbf{Weisfeiler-Lehman (WL) Kernel} \cite{shervashidze2011weisfeiler}: Iteratively relabels nodes based on neighbor signatures (see Section \ref{sec:wl}).
    \item \textbf{Graphlet Kernel} \cite{graphkernels}: Counts small, non-isomorphic subgraphs.
    \item \textbf{Random Walk Kernel} \cite{vishwanathan2010graph}: Compares graphs based on matching paths.
\end{itemize}
WL-based features are widely used in cheminformatics and computer vision \cite{young2022literature, simonovsky2017dynamic}  as well as in neighbour fields like planning \cite{chen2024return}. Furthermore, the message-passing mechanism in Graph Neural Networks is fundamentally inspired by the WL algorithm \cite{Xuetal2018, morris2023weisfeiler}.